\chapter{What was wrong in the legacy simulator}
\label{ch:audit}

The audit ran the legacy simulator on the real SeQUeNCe stack with an execution-event recorder and an invariant checker
attached (Chapter~\ref{ch:corrections}). Each finding below states its evidence: \emph{source} (code read), \emph{measured}
(run), or both.

\begin{table}[H]
\centering\small
\begin{tabular}{@{}p{0.05\linewidth}p{0.36\linewidth}p{0.52\linewidth}@{}}
\toprule
 & Question & Answer \\ \midrule
F1 & Does ACE let two requests share a core? & No. \code{serialize_core_conflicts} splits every layer into batches with at most one request per core, whenever a compiler controller is attached. On QFT the peak demand per core falls from 24 to 1, so the memory budget cannot bind. Source and measured. \\
F2 & Can two requests be mistaken for one? & \textbf{Yes.} SeQUeNCe's \code{Reservation} equality compares initiator, responder, window, size and fidelity, not the request identity. Serialization hid this in every legacy run. Source and measured. \\
F3 & Does a rejection release earlier bookings? & Yes, exactly once on every node it travelled back through. \\
F4 & Does the physics layer limit a link to one pair at a time? & \textbf{No.} 1 to 16 simultaneous requests on one link completed with zero retries. \\
F5 & Phantom occupancy from compiler preparations? & Not in the 24 legacy audits: every booking was released. \\
F6 & Is a trace layer a real barrier? & A conservative one; it is kept. \\
F7 & A request that can never be admitted? & It retries for ever (found with a mistaken test case). \\
F8 & Is waiting measured faithfully? & No: admission was polled every 1\,ms. \\
F9 & Resubmission at the instant a booking expires? & With short windows, stale or double bookings appear in 4 of 6 seeds; a delay of 1\,ps removes them. \\
F10 & Polling admission with one memory per core? & Two requests that swap two cores livelock: 25\,000 submissions each. \\
\bottomrule
\end{tabular}
\caption{Findings of the audit. Full evidence is in \code{docs/reports/ace-correctness-audit.md}.}
\end{table}

\section{Reservation aliasing (F2)}
With four requests on one link and 32 memories per core, the legacy code gave latencies of $1.2/55.7/62.2/67.9$\,ms with three
retries and two orphaned pairs; the same requests, once they are distinct bookings, take $2.4/0.4/1.6/1.2$\,ms with no retry. The
mechanism is that releasing a booking removes the \emph{first equal} reservation, rules are matched to \emph{any equal}
reservation, and the application's dictionaries are keyed by reservation. Reverting only the equality reproduces the old
behaviour and the checker flags it, which shows that the fix is the cause.

\section{Event order is undefined (F9)}
SeQUeNCe's event list is a binary heap ordered by time and priority only, so events at the same instant have no defined order. The
legacy manager started a retry at the very instant the old booking expired. In an A/B test the only change was to resubmit 1\,ps later:
the violations disappeared. No code in the corrected scheduler relies on one same-instant event preceding another.

\section{Livelock under polling (F10)}
Randomized small traces found a request set that never completes under polling admission: requests 1 (core 8 to 12) and 3 (12 to 8),
with one memory per core, each book their own memory, find the other's taken, roll back and retry on the same 1\,ms period. The
invariants held (nothing was corrupted), so this is a livelock, not corruption.

\section{What the audit did not cover}
One mapping seed and one physics seed for the 24 legacy audits; the autonomous strategies; and expiry of unused cached pairs, which
was found later to never occur in the default runs and was then forced (Section~\ref{sec:expiry}).
